% ============================================================
%  Criterios de éxito y plan de validación
% ============================================================
\section{Criterios de evaluación y éxito}
\label{sec:criterios_evaluacion}

El sistema se considera exitoso si cumple con los siguientes criterios:
\begin{itemize}
    \item \textbf{Tiempo de prueba:} reducción de al menos \SI{30}{\percent} del
          tiempo de ensayo respecto de los métodos manuales, gracias a la
          secuencia automática y a la eliminación de ajustes manuales.
    \item \textbf{Automatización:} reducción de al menos \SI{80}{\percent} de la
          intervención del operario en configuración, conmutación entre ensayos y
          adquisición de datos.
    \item \textbf{Seguridad:} alarmas ante ensayos no aprobados o exposición a
          tensiones peligrosas, y comportamiento seguro ante fallas.
    \item \textbf{Escalabilidad:} adaptación a nuevos tipos de dispositivos y a
          cambios normativos sin reestructurar el sistema.
    \item \textbf{Cumplimiento normativo:} ensayos realizados conforme a la
          IEC~60335 y mediciones evaluadas contra sus límites, con trazabilidad.
    \item \textbf{Cronograma:} respeto de las fases y plazos planificados
          (Anexo~\ref{anx:cronograma}), con una tolerancia máxima de diez días.
    \item \textbf{Presupuesto:} costo de desarrollo dentro de \SI{1030}{USD},
          con un margen razonable solo ante contingencias justificadas.
\end{itemize}

\section{Plan de validación}
\label{sec:plan_validacion}

La validación se plantea en tres niveles complementarios: hardware, software y
sistema integrado.

\subsection{Validación del hardware}
\begin{itemize}
    \item \textbf{Verificación funcional:} cada módulo (alimentación, protección,
          medición y desconexión) se evalúa por separado frente a sus
          especificaciones.
    \item \textbf{Pruebas de seguridad:} aislamiento, funcionamiento de las
          protecciones y activación de alarmas ante fallas.
    \item \textbf{Calibración:} contraste con instrumentos patrón certificados,
          cuya incertidumbre sea al menos diez veces menor que la del
          dispositivo.
\end{itemize}

\subsection{Validación del software}
\begin{itemize}
    \item \textbf{Pruebas unitarias:} veredictos, ensayos con hardware simulado,
          relés, base de datos y pestañas de la interfaz.
    \item \textbf{Pruebas de integración:} comunicación entre módulos y con el
          hardware, y correcta transferencia y almacenamiento de datos.
    \item \textbf{Usabilidad y seguridad:} configuración y ejecución intuitivas,
          con mensajes claros de error o advertencia.
\end{itemize}

\subsection{Validación del sistema frente a la norma IEC~60335}
\begin{itemize}
    \item \textbf{Ensayos comparativos:} mediciones con el prototipo y con un
          instrumento calibrado comercial, para verificar la concordancia dentro
          del margen normativo.
    \item \textbf{Verificación con patrones:} el sistema incluye un modo en el que
          se ensayan dos patrones conocidos; el primero debe aprobar y el segundo
          fallar en fuga, puesta a tierra y rigidez, y cada medición debe quedar
          dentro de \SI{\pm 10}{\percent} del valor del patrón.
    \item \textbf{Intercomparación:} cuando sea posible, contraste con resultados
          emitidos por un laboratorio acreditado.
    \item \textbf{Criterio de aceptación:} el sistema es válido si emite el
          veredicto correcto frente a los límites de la norma
          (Anexo~\ref{anx:valores_referencia}) con trazabilidad demostrable.
\end{itemize}
